\paragrafos{Πειράματα τύχης και δειγματικός χώρος}{Γ.3.1, «Πειράματα τύχης, δειγματικός χώρος»}

\begin{anagke}
Ο Νίκος και η Ρήνη κληρώνουν με ένα νόμισμα ποιος θα ξεκινήσει το παιχνίδι. Ο Νίκος έχασε τρεις φορές στη σειρά: έφερε «Κεφάλι» και τις τρεις. Την τέταρτη φορά λέει με σιγουριά: «Τώρα θα φέρει Γράμματα, αποκλείεται να ξαναφέρει Κεφάλι!». Η Ρήνη του απαντά: «Το νόμισμα δεν θυμάται τι έγινε πριν».

Η κυρία Αλεξάνδρα ρωτά: «Τι μπορείτε να πείτε με βεβαιότητα πριν ρίξει ο Νίκος;». Η Φανή σκέφτεται λίγο: «Ότι θα είναι είτε Κεφάλι είτε Γράμματα. Τίποτε άλλο δεν γίνεται». Αυτό ακριβώς είναι το πρώτο βήμα στην επιστήμη της τύχης: να μη μαντεύουμε, αλλά να γράψουμε πρώτα όλα όσα μπορούν να συμβούν.

Σε αυτή την παράγραφο μαθαίνουμε ποια πειράματα λέγονται πειράματα τύχης και πώς γράφουμε το σύνολο όλων των δυνατών αποτελεσμάτων τους, που το λέμε δειγματικό χώρο.
\end{anagke}

\begin{lexi}{πείραμα τύχης, δειγματικός χώρος}{ονόματα και σημασίες}
Το \lexiword{πείραμα} είναι μια δοκιμή. Η \lexiword{τύχη} είναι ό,τι δεν μπορούμε να προβλέψουμε. Το \lexiword{δείγμα} (από το «δείχνω») είναι ένα παράδειγμα, ένα αποτέλεσμα που βλέπουμε. Ο \lexiword{δειγματικός χώρος} είναι ο «χώρος» που χωράει όλα τα δείγματα, δηλαδή όλα τα δυνατά αποτελέσματα.
\end{lexi}

\begin{metafrasi}
Ρίχνω ένα ζάρι και δεν ξέρω τι θα φέρει. & πείραμα τύχης\\
Ένα από τα πράγματα που μπορούν να συμβούν. & αποτέλεσμα\\
Όλα όσα μπορούν να συμβούν. & δειγματικός χώρος\\
Οι δυνατές έδρες ενός ζαριού. & $\{1,2,3,4,5,6\}$\\
Πρώτα Κεφάλι και μετά Γράμματα. & $\mathrm{ΚΓ}$\\
\end{metafrasi}

\tmima{Η θεωρία}

\begin{orismos}{πείραμα τύχης}
\textbf{Πείραμα τύχης} λέγεται ένα πείραμα στο οποίο \emph{δεν μπορούμε να προβλέψουμε με βεβαιότητα το αποτέλεσμα}, όσες φορές κι αν το επαναλάβουμε κάτω από τις ίδιες συνθήκες.
\end{orismos}

Η ρίψη ενός ζαριού και η ρίψη ενός νομίσματος είναι πειράματα τύχης. Η κλήρωση ενός ονόματος από μια κάλπη και η περιστροφή ενός τροχού στο λούνα παρκ επίσης. Αντίθετα, το πείραμα «βράζω νερό στους $100$ $^\circ$C σε κανονική πίεση και κοιτάζω αν βράζει» δεν είναι πείραμα τύχης, γιατί ξέρουμε το αποτέλεσμα πριν το κάνουμε.

\begin{orismos}{δειγματικός χώρος}
Κάθε δυνατό αποτέλεσμα ενός πειράματος τύχης λέγεται \textbf{στοιχείο} του δειγματικού χώρου. Το σύνολο όλων των δυνατών αποτελεσμάτων λέγεται \textbf{δειγματικός χώρος} του πειράματος. Τον γράφουμε μέσα σε άγκιστρα, χωρίζοντας τα στοιχεία με κόμμα. Το \textbf{πλήθος} των στοιχείων του είναι το πλήθος των δυνατών αποτελεσμάτων.
\end{orismos}

\begin{center}
\bindie1\hspace{0.5em}\bindie2\hspace{0.5em}\bindie3\hspace{0.5em}\bindie4\hspace{0.5em}\bindie5\hspace{0.5em}\bindie6\\[1ex]
{\small Δειγματικός χώρος ζαριού: $\{1,2,3,4,5,6\}$, έξι στοιχεία}
\end{center}

Για τη ρίψη ενός νομίσματος ο δειγματικός χώρος είναι $\{\mathrm{Κ},\ \mathrm{Γ}\}$ (Κεφάλι, Γράμματα), δύο στοιχεία.

\begin{kanonas}{πώς γράφω τον δειγματικό χώρο}
1. Καταλαβαίνω τι ακριβώς καταγράφω σε κάθε εκτέλεση του πειράματος.\\
2. Γράφω τα δυνατά αποτελέσματα με σύστημα (π.χ. με σειρά: πρώτα όσα ξεκινούν από το πρώτο αποτέλεσμα κ.ο.κ.).\\
3. Ελέγχω ότι κανένα δεν λείπει και κανένα δεν γράφτηκε δύο φορές.
\end{kanonas}

Όταν το πείραμα έχει δύο ή περισσότερα στάδια, π.χ. ρίχνουμε δύο φορές το νόμισμα, το αποτέλεσμα είναι \emph{διατεταγμένο}: καταγράφουμε τι έγινε πρώτα και τι έγινε μετά. Ο δειγματικός χώρος είναι $\{\mathrm{ΚΚ},\ \mathrm{ΚΓ},\ \mathrm{ΓΚ},\ \mathrm{ΓΓ}\}$ και έχει τέσσερα στοιχεία: το $\mathrm{ΚΓ}$ και το $\mathrm{ΓΚ}$ είναι διαφορετικά αποτελέσματα.

\begin{prosoxi}
\begin{itemize}[leftmargin=1.2em,itemsep=0.3ex,topsep=2pt]
\item Ο δειγματικός χώρος εξαρτάται από \emph{το τι καταγράφουμε}. Αν ρίξουμε ζάρι και γράψουμε «άρτιος ή περιττός», ο δειγματικός χώρος είναι $\{\text{άρτιος},\ \text{περιττός}\}$ και όχι $\{1,\dots,6\}$.
\item Στα πειράματα με δύο στάδια το $\mathrm{ΚΓ}$ και το $\mathrm{ΓΚ}$ είναι \emph{διαφορετικά} αποτελέσματα.
\item Στον δειγματικό χώρο γράφουμε όλα τα δυνατά αποτελέσματα, ακόμη και αυτά που μας φαίνονται λιγότερο πιθανά.
\item Το ότι βγήκε τρεις φορές Κεφάλι δεν αλλάζει τον δειγματικό χώρο: την επόμενη φορά μπορεί να έρθει και πάλι Κ ή Γ.
\end{itemize}
\end{prosoxi}

\tmima{Λυμένα παραδείγματα}

\paradeigma{Είναι πείραμα τύχης;}
\tixero Τέσσερις καταστάσεις: α) Ρίχνω ένα ζάρι και βλέπω την πάνω έδρα. β) Υπολογίζω το $8\cdot7$. γ) Παίρνω ένα φύλλο από μια ανακατεμένη τράπουλα. δ) Βλέπω τι μέρα της εβδομάδας είναι αύριο, αν σήμερα είναι Τρίτη.\par
\tithelo Να ξεχωρίσω ποιες είναι πειράματα τύχης.\par
\possindeo Ρωτάω αν μπορώ να προβλέψω το αποτέλεσμα με βεβαιότητα.\par
\apantisi Πειράματα τύχης είναι τα α) και γ). Το β) έχει πάντα αποτέλεσμα $56$ και το δ) είναι Τετάρτη, άρα και τα δύο προβλέπονται.\par
\elegxos Στα α) και γ) ούτε με πολλές επαναλήψεις θα έβρισκα ποιο αποτέλεσμα θα έρθει την επόμενη φορά.

\paradeigma{Δύο ρίψεις νομίσματος}
\tixero Ρίχνουμε ένα νόμισμα δύο φορές και σημειώνουμε με τη σειρά τι έφερε.\par
\tithelo Τον δειγματικό χώρο.\par
\possindeo Γράφω πρώτα όσα ξεκινούν από Κ και μετά όσα ξεκινούν από Γ.\par
\apantisi $\{\mathrm{ΚΚ},\ \mathrm{ΚΓ},\ \mathrm{ΓΚ},\ \mathrm{ΓΓ}\}$, με $4$ στοιχεία.\par
\elegxos Για κάθε πρώτο αποτέλεσμα ($2$ επιλογές) υπάρχουν δύο δεύτερα, άρα $2\cdot2=4$.

\paradeigma{Με επανατοποθέτηση και χωρίς}
\tixero Ένα κουτί έχει τρεις μπάλες: κόκκινη (Κ), μπλε (Μ) και πράσινη (Π). Βγάζουμε δύο μπάλες, τη μία μετά την άλλη, και σημειώνουμε τα χρώματα με τη σειρά.\par
\tithelo Τον δειγματικό χώρο α) αν δεν ξαναβάζουμε την πρώτη μπάλα στο κουτί, β) αν την ξαναβάζουμε.\par
\possindeo Ταξινομώ τα αποτελέσματα κατά την πρώτη μπάλα.\par
\apantisi α) $\{\mathrm{ΚΜ},\ \mathrm{ΚΠ},\ \mathrm{ΜΚ},\ \mathrm{ΜΠ},\ \mathrm{ΠΚ},\ \mathrm{ΠΜ}\}$, $6$ στοιχεία. β) Τα $6$ του α) και ακόμη τα $\mathrm{ΚΚ},\ \mathrm{ΜΜ},\ \mathrm{ΠΠ}$: συνολικά $9$ στοιχεία.\par
\elegxos Στο α) η πρώτη μπάλα έχει $3$ επιλογές και η δεύτερη $2$ ($3\cdot2=6$). Στο β) κάθε φορά έχω $3$ επιλογές ($3\cdot3=9$).

\paradeigma{Ο δειγματικός χώρος εξαρτάται από το τι καταγράφω}
\tixero Ρίχνουμε ένα ζάρι.\par
\tithelo Τον δειγματικό χώρο α) αν καταγράφουμε την πάνω έδρα, β) αν καταγράφουμε αν ο αριθμός είναι άρτιος ή περιττός, γ) αν καταγράφουμε αν έφερε $6$ ή όχι.\par
\possindeo Για κάθε περίπτωση σκέφτομαι ποιες διαφορετικές «απαντήσεις» μπορεί να πάρω.\par
\apantisi α) $\{1,2,3,4,5,6\}$ ($6$ στοιχεία). β) $\{\text{άρτιος},\ \text{περιττός}\}$ ($2$ στοιχεία). γ) $\{6,\ \text{όχι }6\}$ ($2$ στοιχεία).\par
\elegxos Το ίδιο φυσικό πείραμα έδωσε τρεις διαφορετικούς δειγματικούς χώρους, γιατί αλλάξαμε το τι σημειώνουμε.

\tmima{Ασκήσεις}

\zoni{Το κατάλαβα}{πειράματα τύχης και δειγματικοί χώροι}

\ex Ποια από τα παρακάτω είναι πειράματα τύχης; α) Ρίχνω ζάρι και καταγράφω την πάνω έδρα. β) Προσθέτω $17+25$. γ) Γυρίζω τον τροχό του λούνα παρκ και βλέπω πού σταματά. δ) Κοιτάζω το ρολόι και γράφω τι ώρα είναι. ε) Κληρώνω ένα όνομα από τον κατάλογο της τάξης. στ) Βάζω νερό στους $-5$ $^\circ$C και βλέπω αν παγώνει.

\ex Γράψε τον δειγματικό χώρο: α) Ρίχνω ένα νόμισμα. β) Ρίχνω ένα ζάρι. γ) Παίρνω μια μπάλα από ένα κουτί με μια άσπρη, μια μαύρη και μια κόκκινη μπάλα. δ) Ρίχνω ζάρι και σημειώνω αν ο αριθμός είναι άρτιος ή περιττός.

\ex Γράψε τον δειγματικό χώρο και πες πόσα στοιχεία έχει: α) Επιλέγω τυχαία μια μέρα της εβδομάδας με όνομα που αρχίζει από Τ. β) Επιλέγω τυχαία ένα ψηφίο του αριθμού $2024$. γ) Επιλέγω τυχαία ένα γράμμα της λέξης ΠΙΘΑΝΟΤΗΤΑ.

\ex Ρίχνουμε ένα ζάρι και καταγράφουμε: α) αν έφερε $6$ ή όχι, β) το υπόλοιπο της διαίρεσης του αριθμού με το $3$, γ) αν ο αριθμός είναι πρώτος ή όχι. Βρες τον δειγματικό χώρο σε κάθε περίπτωση.

\ex Βρες ένα πείραμα τύχης που να έχει δειγματικό χώρο: α) $\{1,2,3,4,5,6,7,8\}$. β) $\{\text{ΝΑΙ},\ \text{ΟΧΙ}\}$. γ) $\{\text{Δευτέρα},\ \text{Τρίτη},\ \text{Τετάρτη}\}$.

\ex Σωστό ή λάθος; α) Σε ένα πείραμα τύχης ξέρουμε πάντα το αποτέλεσμα πριν το κάνουμε. β) Ο δειγματικός χώρος περιέχει όλα τα δυνατά αποτελέσματα. γ) Στις δύο ρίψεις νομίσματος το $\mathrm{ΚΓ}$ και το $\mathrm{ΓΚ}$ είναι το ίδιο αποτέλεσμα. δ) Ο δειγματικός χώρος της ρίψης ενός ζαριού έχει $6$ στοιχεία. ε) Ένα αποτέλεσμα που είναι πολύ δύσκολο να συμβεί δεν γράφεται στον δειγματικό χώρο.

\zoni{Το χρησιμοποιώ}{κουτιά, αγώνες, κωδικοί}

\ex \textbf{Οι καραμέλες.} Ένα κουτί έχει δύο καραμέλες: φράουλα (Φ) και λεμόνι (Λ). Παίρνουμε δύο καραμέλες την μία μετά την άλλη. Βρες τον δειγματικό χώρο α) αν δεν ξαναβάζουμε την πρώτη, β) αν την ξαναβάζουμε.

\ex \textbf{Δύο αγώνες.} Η ομάδα της Α3 παίζει δύο αγώνες. Κάθε αγώνας τελειώνει με νίκη (Ν), ισοπαλία (Ι) ή ήττα (Η). Γράψε τον δειγματικό χώρο των δύο αγώνων και πες πόσα στοιχεία έχει.

\ex \textbf{Το τηλέφωνο.} Η Φανή θυμάται το τηλέφωνο της θείας Μαρίας, εκτός από το τελευταίο ψηφίο. Ξέρει ότι είναι άρτιος αριθμός. Γράψε τον δειγματικό χώρο για το ψηφίο που λείπει.

\ex \textbf{Οι διψήφιοι.} Με τα ψηφία $3$ και $7$ γράφουμε διψήφιους αριθμούς (επιτρέπεται να επαναληφθεί το ίδιο ψηφίο). Γράψε τον δειγματικό χώρο.

\ex \textbf{Η ντουλάπα.} Η Αγγελική διαλέγει ένα παντελόνι από δύο (τζιν Τ, φόρμα Φ) και μια μπλούζα από δύο (άσπρη Α, κίτρινη Κ). Γράψε τον δειγματικό χώρο του πειράματος να διαλέξει ένα από κάθε είδος (πρώτα παντελόνι, μετά μπλούζα).

\zoni{Το σκέφτομαι}{όλα τα δυνατά και ο τρόπος καταγραφής}

\ex Ρίχνουμε ένα νόμισμα τρεις φορές και γράφουμε με τη σειρά τι έφερε. Γράψε τον δειγματικό χώρο και πες πόσα στοιχεία έχει.

\ex Ο Νίκος λέει: «Έφερα Κεφάλι τρεις φορές, άρα την επόμενη θα φέρω σίγουρα Γράμματα». Έχει δίκιο; Εξήγησε με βάση τον δειγματικό χώρο.

\ex Ο Στέλιος ρίχνει δύο ζάρια και λέει: «Ο δειγματικός χώρος του ενός ζαριού έχει $6$ στοιχεία, άρα των δύο ζαριών $6+6=12$». Βρες το σωστό πλήθος, γράφοντας για κάθε αποτέλεσμα του πρώτου ζαριού τα δυνατά αποτελέσματα του δεύτερου.

\ex Μπορεί ένα πείραμα τύχης να έχει δειγματικό χώρο με ένα μόνο στοιχείο; Δώσε παράδειγμα ή εξήγησε γιατί όχι. Αν πάλι ναι, είναι τότε «τύχη» το αποτέλεσμα;

\zoni{Για το σπίτι}{πειράματα γύρω μας}

\ex Βρες τρία πειράματα τύχης από την καθημερινή σου ζωή και γράψε για το καθένα τον δειγματικό χώρο.

\ex Μια τράπουλα έχει $52$ φύλλα σε τέσσερις σειρές: κούπες, καρό, σπαθιά και μπαστούνια, με $13$ φύλλα στην καθεμία. Τραβάμε ένα φύλλο και σημειώνουμε μόνο τη σειρά του. Γράψε τον δειγματικό χώρο. Πόσα στοιχεία θα είχε αν σημειώναμε ολόκληρο το φύλλο;
